\chapter{Arquitectura de control y navegación: \texttt{ros2\_control}, SLAM y Nav2}
\label{chap:app_control_nav}

Este capítulo describe la arquitectura de software que integra control diferencial, mapeo, navegación autónoma y exploración de fronteras. Las capas de hardware, firmware y gemelo digital fueron desarrolladas en capítulos anteriores; por ello, aquí se enfatiza la orquestación de los subsistemas mediante \texttt{ros2\_control}, \texttt{slam\_toolbox}, Nav2 y \texttt{m-explore-ros2}. Los parámetros completos de configuración se trasladan al \ref{anx:nav2_slam_config}.

%------------------------------------------------------------------
\section{Marco de control: \texttt{ros2\_control}}
\label{sec:ros2ctrl_overview}

El marco \texttt{ros2\_control} permite usar una misma interfaz de control tanto en simulación como en hardware real. El \texttt{controller\_manager} gestiona el \texttt{diff\_drive\_controller}, que convierte comandos de velocidad $(v,\omega)$ en velocidades de rueda y publica la odometría, y el \texttt{joint\_state\_broadcaster}, encargado de actualizar el estado de las articulaciones.

En hardware real, la interfaz se conecta al \gls{Arduino} mediante el plugin \texttt{DiffDriveArduino}; en simulación, Gazebo reemplaza al hardware físico. En ambos casos se conservan los mismos marcos, tópicos y parámetros geométricos, facilitando la transferencia simulación--hardware.

%------------------------------------------------------------------
\section{Mapeo con \texttt{slam\_toolbox}}
\label{sec:slam_ch9}

El robot construye el mapa del entorno mediante \texttt{slam\_toolbox}, que estima simultáneamente la trayectoria del robot y el mapa 2D a partir de la odometría y los barridos del \gls{LIDAR}. Se utiliza el modo \texttt{online\_async}, ya que procesa los datos del láser sin bloquear la adquisición de nuevos barridos, manteniendo una cadena \texttt{tf} estable en una Raspberry Pi 4.

El mapeo utiliza el marco \texttt{odom} como referencia odométrica, \texttt{map} como marco global, \texttt{base\_link} como base del robot y \texttt{/scan} como tópico del sensor láser. La configuración activa cierre de lazo y optimización de grafo mediante Ceres, permitiendo corregir deriva acumulada durante el recorrido.

\begin{table}[H]
\centering
\caption{Parámetros principales de \texttt{slam\_toolbox}.}
\label{tab:slam_params_ch9}
\small
\begin{tabular}{lll}
\toprule
\textbf{Parámetro} & \textbf{Valor} & \textbf{Función} \\
\midrule
\texttt{solver\_plugin} & \texttt{CeresSolver} & Optimización del grafo \\
\texttt{odom\_frame} & \texttt{odom} & Marco de odometría \\
\texttt{map\_frame} & \texttt{map} & Marco global \\
\texttt{base\_frame} & \texttt{base\_link} & Marco base del robot \\
\texttt{scan\_topic} & \texttt{/scan} & Entrada del \gls{LIDAR} \\
\texttt{resolution} & 0.05~m & Resolución del mapa \\
\texttt{do\_loop\_closing} & \texttt{true} & Cierre de lazo activado \\
\bottomrule
\end{tabular}
\end{table}

%------------------------------------------------------------------
\section{Navegación autónoma con Nav2}
\label{sec:nav2_ch9}

Nav2 recibe una meta de navegación, calcula una ruta global y genera comandos de velocidad para que el robot la ejecute evitando obstáculos. La cadena funcional implementada es:

\begin{center}
\small
Meta $\rightarrow$ \texttt{bt\_navigator} $\rightarrow$ \texttt{planner\_server} $\rightarrow$ \texttt{controller\_server}
$\rightarrow$ \texttt{twist\_mux} $\rightarrow$ \texttt{diff\_drive\_controller}
\end{center}

El planificador global utilizado es NavFn con A\*, activado mediante \texttt{use\_astar}. Este planificador opera sobre el costmap global y genera una trayectoria desde la pose actual hasta la meta. Para el seguimiento local se emplea DWB, que evalúa trayectorias candidatas y selecciona comandos de velocidad compatibles con las restricciones cinemáticas del robot.

\begin{table}[H]
\centering
\caption{Componentes principales de Nav2 utilizados.}
\label{tab:nav2_componentes_ch9}
\small
\begin{tabular}{lll}
\toprule
\textbf{Componente} & \textbf{Implementación} & \textbf{Función} \\
\midrule
Planificador global & NavFn con A\* & Generar ruta global \\
Controlador local & DWB & Seguir ruta y evitar obstáculos \\
Costmaps & Global y local & Representar obstáculos y zonas transitables \\
Árbol de comportamiento & \texttt{bt\_navigator} & Orquestar navegación y recuperación \\
Multiplexor & \texttt{twist\_mux} & Priorizar navegación o teleoperación \\
\bottomrule
\end{tabular}
\end{table}

%------------------------------------------------------------------
\section{Mapas de costos}
\label{sec:costmaps_ch9}

Nav2 utiliza dos mapas de costos. El \texttt{global\_costmap} entrega la representación usada por el planificador global, mientras que el \texttt{local\_costmap} se actualiza alrededor del robot y permite al controlador DWB reaccionar ante obstáculos cercanos. Ambos mapas incorporan una capa de inflación, necesaria para mantener una distancia de seguridad respecto de los obstáculos.

En hardware real se utiliza una resolución de 0.05~m/celda y un radio de inflación de 0.15~m. En simulación se emplean parámetros más permisivos, adecuados al modelo virtual del robot. El detalle completo de frecuencias, radios y pesos se presenta en el \ref{anx:nav2_slam_config}.

%------------------------------------------------------------------
\section{Árbol de comportamiento y recuperación}
\label{sec:bt_ch9}

La navegación de alto nivel se organiza mediante un árbol de comportamiento. La rama principal calcula la ruta global y ejecuta el seguimiento local; si ocurre un fallo, se activan acciones de recuperación. Esta estructura permite reintentos y recuperación ante bloqueos sin modificar la lógica completa del sistema.

Las acciones de recuperación configuradas incluyen limpieza de costmaps, retroceso, giro en el lugar, espera y replanificación. La lógica sigue un principio de escalada: primero acciones de bajo costo computacional y luego maniobras físicas, como retroceder o girar.

\begin{table}[H]
\centering
\caption{Acciones de recuperación consideradas en Nav2.}
\label{tab:recovery_ch9}
\small
\begin{tabular}{ll}
\toprule
\textbf{Acción} & \textbf{Propósito} \\
\midrule
\textit{Clear local costmap} & Eliminar obstáculos fantasma cercanos \\
\textit{Clear global costmap} & Limpiar obstáculos espurios del mapa global \\
\textit{BackUp} & Retroceder para liberar espacio \\
\textit{Spin} & Girar para obtener nueva perspectiva \\
\textit{Wait} & Esperar antes de reintentar \\
\bottomrule
\end{tabular}
\end{table}

%------------------------------------------------------------------
\section{Exploración autónoma de fronteras}
\label{sec:explore_ch9}

Para mapear entornos desconocidos sin intervención manual continua, se integra un módulo de exploración basado en fronteras. Se utiliza \texttt{m-explore-ros2}, derivado de \texttt{explore\_lite}, cuyo nodo \texttt{explore\_server} identifica fronteras entre espacio libre conocido y zonas no observadas, y envía metas de navegación a Nav2.

El explorador consume el costmap global, evalúa las fronteras disponibles y selecciona metas considerando distancia y ganancia de información. De esta forma, se cierra el ciclo percepción--mapeo--navegación--exploración sobre la misma arquitectura.

\begin{table}[H]
\centering
\caption{Parámetros principales del explorador de fronteras.}
\label{tab:explore_params_ch9}
\small
\begin{tabular}{lll}
\toprule
\textbf{Parámetro} & \textbf{Valor} & \textbf{Función} \\
\midrule
\texttt{robot\_base\_frame} & \texttt{base\_link} & Marco base del robot \\
\texttt{costmap\_topic} & \texttt{/global\_costmap/costmap} & Costmap usado para detectar fronteras \\
\texttt{planner\_frequency} & 0.05~Hz & Frecuencia de reevaluación \\
\texttt{progress\_timeout} & 60.0~s & Tiempo máximo sin progreso \\
\texttt{gain\_scale} & 1.0 & Peso de ganancia de información \\
\texttt{potential\_scale} & 3.0 & Peso de distancia a la frontera \\
\bottomrule
\end{tabular}
\end{table}

%------------------------------------------------------------------
\section{Integración del ciclo autónomo}
\label{sec:integracion_ciclo}

Los subsistemas operan de forma encadenada. \texttt{slam\_toolbox} construye el mapa y publica la transformación \texttt{map}$\rightarrow$\texttt{odom}; Nav2 utiliza ese mapa para planificar y controlar el movimiento; \texttt{twist\_mux} envía los comandos al controlador diferencial; y \texttt{explore\_server} genera nuevas metas hacia fronteras no exploradas.

Esta integración permite que el robot pase desde teleoperación y mapeo inicial hacia navegación autónoma y exploración de zonas desconocidas, manteniendo la misma base de marcos y tópicos tanto en simulación como en hardware.

%------------------------------------------------------------------
\section{Transferencia simulación--hardware}
\label{sec:verificacion_ch9}

La transferencia al hardware se realiza manteniendo la estructura de tópicos y marcos, pero seleccionando archivos de configuración específicos para el entorno real. En simulación se emplean archivos con \texttt{use\_sim\_time: true}; en hardware, los archivos equivalentes usan \texttt{use\_sim\_time: false}. Además, se ajustan límites cinemáticos, frecuencias, tolerancias e inflación de costmaps según las capacidades del prototipo físico.

%------------------------------------------------------------------
\section*{Conclusión del capítulo}
\label{sec:conclusion_ch9}

Este capítulo describió la arquitectura de navegación autónoma del robot, integrando \texttt{ros2\_control}, \texttt{slam\_toolbox}, Nav2, mapas de costos, árbol de comportamiento y exploración basada en fronteras. La arquitectura permite articular percepción, mapeo, navegación y exploración sobre una misma base de tópicos y marcos, tanto en simulación como en hardware real. Su desempeño se evalúa en el Capítulo~\ref{ch:resultados}.